\documentclass[11pt]{article}

% ===================================================================
% Shared preamble for the "tiny lemma, read slowly" preprint series.
% Mirrors the style of FrobeniusLemmas_preprint.tex.
% ===================================================================

\usepackage{amsmath,amssymb,amsthm}
\usepackage{mathtools}
\usepackage[margin=1.2in]{geometry}
\usepackage{hyperref}
\usepackage{xcolor}
\usepackage{listings}
\usepackage{tikz}
\usepackage{tikz-cd}
\usepackage{mathpartir}
\usepackage{newunicodechar}
\usepackage{setspace}

\onehalfspacing

% ---- Unicode glyphs ----
\newunicodechar{→}{\ensuremath{\to}}
\newunicodechar{↦}{\ensuremath{\mapsto}}
\newunicodechar{↔}{\ensuremath{\leftrightarrow}}
\newunicodechar{⟹}{\ensuremath{\Longrightarrow}}
\newunicodechar{⟶}{\ensuremath{\longrightarrow}}
\newunicodechar{≠}{\ensuremath{\neq}}
\newunicodechar{≤}{\ensuremath{\leq}}
\newunicodechar{≥}{\ensuremath{\geq}}
\newunicodechar{∀}{\ensuremath{\forall}}
\newunicodechar{∃}{\ensuremath{\exists}}
\newunicodechar{∘}{\ensuremath{\circ}}
\newunicodechar{·}{\ensuremath{\cdot}}
\newunicodechar{∈}{\ensuremath{\in}}
\newunicodechar{∉}{\ensuremath{\notin}}
\newunicodechar{≃}{\ensuremath{\simeq}}
\newunicodechar{≅}{\ensuremath{\cong}}
\newunicodechar{≈}{\ensuremath{\approx}}
\newunicodechar{∧}{\ensuremath{\wedge}}
\newunicodechar{∨}{\ensuremath{\vee}}
\newunicodechar{¬}{\ensuremath{\neg}}
\newunicodechar{⊢}{\ensuremath{\vdash}}
\newunicodechar{⟨}{\ensuremath{\langle}}
\newunicodechar{⟩}{\ensuremath{\rangle}}
\newunicodechar{⇑}{\ensuremath{\mathord{\uparrow}}}
\newunicodechar{↑}{\ensuremath{\mathord{\uparrow}}}
\newunicodechar{∎}{\ensuremath{\blacksquare}}
\newunicodechar{∣}{\ensuremath{\mid}}
\newunicodechar{∤}{\ensuremath{\nmid}}
\newunicodechar{∑}{\ensuremath{\textstyle\sum}}
\newunicodechar{∏}{\ensuremath{\textstyle\prod}}
\newunicodechar{∅}{\ensuremath{\emptyset}}
\newunicodechar{∪}{\ensuremath{\cup}}
\newunicodechar{∩}{\ensuremath{\cap}}
\newunicodechar{≡}{\ensuremath{\equiv}}
\newunicodechar{ℕ}{\ensuremath{\mathbb{N}}}
\newunicodechar{ℤ}{\ensuremath{\mathbb{Z}}}
\newunicodechar{ℝ}{\ensuremath{\mathbb{R}}}
\newunicodechar{𝔽}{\ensuremath{\mathbb{F}}}
\newunicodechar{α}{\ensuremath{\alpha}}
\newunicodechar{β}{\ensuremath{\beta}}
\newunicodechar{σ}{\ensuremath{\sigma}}
\newunicodechar{τ}{\ensuremath{\tau}}
\newunicodechar{φ}{\ensuremath{\varphi}}
\newunicodechar{ψ}{\ensuremath{\psi}}
\newunicodechar{χ}{\ensuremath{\chi}}
\newunicodechar{Φ}{\ensuremath{\Phi}}
\newunicodechar{Δ}{\ensuremath{\Delta}}
\newunicodechar{₀}{\textsubscript{0}}
\newunicodechar{₁}{\textsubscript{1}}
\newunicodechar{₂}{\textsubscript{2}}
\newunicodechar{ᵏ}{\ensuremath{{}^{k}}}
\newunicodechar{ⁿ}{\ensuremath{{}^{n}}}
\newunicodechar{ʲ}{\ensuremath{{}^{j}}}
\newunicodechar{²}{\ensuremath{{}^{2}}}
\newunicodechar{³}{\ensuremath{{}^{3}}}
\newunicodechar{⁴}{\ensuremath{{}^{4}}}
\newunicodechar{⁻}{\ensuremath{{}^{-}}}
\newunicodechar{¹}{\ensuremath{{}^{1}}}

% ---- Lean listing style ----
\definecolor{leanbg}{RGB}{247,247,250}
\definecolor{leankw}{RGB}{0,90,160}
\definecolor{leancm}{RGB}{120,120,120}
\lstdefinestyle{lean}{
  backgroundcolor=\color{leanbg},
  basicstyle=\ttfamily\small,
  keywordstyle=\color{leankw}\bfseries,
  commentstyle=\color{leancm}\itshape,
  morekeywords={theorem,lemma,def,fun,Prop,Type,variable,where,
    Field,Fintype,Finite,DecidableEq,CommRing,CommSemiring,Ring,CharP,Function,
    Injective,Surjective,Bijective,by,rfl,have,rw,exact,ext,simp,intro,
    iterateFrobenius,RingHom,Commute,convert,apply,using,Nat,Coprime,gcd,
    Odd,Even,IsAPN,IsAB,walsh,Nat,obtain,induction,cases,calc,grind},
  showstringspaces=false,
  columns=fullflexible,
  extendedchars=true,
  frame=single,
  framesep=6pt,
  xleftmargin=4pt,
  aboveskip=12pt,
  belowskip=14pt,
}
\lstset{style=lean}

\newtheorem{lemma}{Lemma}
\newtheorem{theorem}{Theorem}
\newtheorem{corollary}{Corollary}

% boxed "what just happened" note
\newcommand{\note}[1]{%
  \par\smallskip
  \noindent\colorbox{leanbg}{\parbox{\dimexpr\linewidth-2\fboxsep}{\small #1}}%
  \par\smallskip}


\title{\textbf{The Shift in Disguise:\\[2pt]
A Non-Faithful Functor from Fields to Symbolic Dynamics}}
\author{}
\date{}

\begin{document}

\maketitle

\thispagestyle{empty}

\begin{abstract}
\noindent
One element of a field becomes one sequence.\\
The Frobenius map becomes the shift.\\
Coprimality becomes a minimal rotation.

\medskip
\noindent
The functor is real, but it forgets.\\
We say exactly what it keeps, and what it drops.

\medskip
\noindent
Everything is checked in Lean~4 / Mathlib.
\end{abstract}

\vfill

\noindent\textit{Conference paper call --- short, slow-read submission.}

\newpage


\section{The picture}

\bigskip

Take a finite field $\mathbb{F}_{2^{\,n}}$.

\bigskip

Pick one element $v$.

\bigskip

Square it, again and again:
\[
  v,\quad v^{2},\quad v^{4},\quad \dots,\quad v^{2^{\,n-1}}.
\]

\bigskip

That is a sequence.

\bigskip

It repeats with period dividing $n$.

\bigskip

\note{One field element. One periodic sequence.}

\newpage


\section{The shift}

\bigskip

Squaring once moves the whole sequence one step left:
\[
  (v,\; v^{2},\; v^{4},\; \dots)
  \;\xmapsto{\ x\,\mapsto\, x^{2}\ }\;
  (v^{2},\; v^{4},\; v^{8},\; \dots).
\]

\bigskip

That is the \emph{shift} $\sigma$.

\bigskip

So the Frobenius map \textbf{is} the shift.
\[
  x \mapsto x^{2} \quad\longleftrightarrow\quad \sigma .
\]

\bigskip

And the $k$-th power of Frobenius is the $k$-fold shift.
\[
  x \mapsto x^{2^{\,k}} \quad\longleftrightarrow\quad \sigma^{k}.
\]

\bigskip

\note{Algebra on the left. Dynamics on the right.}

\newpage


\section{The encoding, in Lean}

\bigskip

The orbit sequence:

\begin{lstlisting}
def frobeniusOrbitSeq (p : ℕ) (x : F) : ℕ → F :=
  fun k => x ^ (p ^ k)
\end{lstlisting}

\bigskip

One step is one Frobenius:

\begin{lstlisting}
theorem frobeniusOrbitSeq_succ (p : ℕ) (x : F) (k : ℕ) :
    frobeniusOrbitSeq p x (k + 1)
      = (frobeniusOrbitSeq p x k) ^ p
\end{lstlisting}

\bigskip

\note{The recurrence \emph{is} the shift rule.}

\newpage


\section{Cool fact 1 --- counting fixed points}

\bigskip

How many sequences does $\sigma^{k}$ fix?

\bigskip

How many $v$ satisfy $v^{2^{\,k}} = v$?

\bigskip

Same question. Same answer.
\[
  \#\{\, v : v^{2^{\,k}} = v \,\} \;=\; 2^{\gcd(k,\,n)} .
\]

\bigskip

The shift side counts \emph{necklaces}:
\[
  \#\,\mathrm{Fix}(\sigma^{k}) \;=\; 2^{\gcd(k,\,n)} .
\]

\bigskip

\note{One formula, read two ways.}

\newpage


\section{Cool fact 1, in Lean}

\bigskip

The field side:

\begin{lstlisting}
theorem frobenius_fixed_card (n k : ℕ) (hn : 0 < n) (hk : 0 < k)
    (hcard : Fintype.card F = 2 ^ n) :
    Nat.card {x : F | x ^ (2 ^ k) = x} = 2 ^ Nat.gcd k n
\end{lstlisting}

\bigskip

The shift side:

\begin{lstlisting}
theorem necklace_orbit_count (n d : ℕ) [NeZero n] (hn : 0 < n) :
    Fintype.card {w : ZMod n → ZMod 2 |
        ∀ i, w i = w (i + (d : ZMod n))}
      = 2 ^ Nat.gcd d n
\end{lstlisting}

\bigskip

\note{Same right-hand side. No coincidence.}

\newpage


\section{Cool fact 2 --- minimal rotations}

\bigskip

Now ask: when is $\sigma^{k}$ a \emph{minimal} rotation of $\mathbb{Z}/n\mathbb{Z}$?

\bigskip

Exactly when
\[
  \gcd(k,\,n) = 1 .
\]

\bigskip

A minimal rotation sweeps the whole cycle.

\bigskip

So it can fix a sequence \textbf{only if the sequence is constant}.

\bigskip

Constant sequence $\;\Longleftrightarrow\;$ the element lives in $\mathbb{F}_{2} = \{0,1\}$.

\bigskip

\note{No invariant sub-coordinate. Nowhere to hide.}

\newpage


\section{Cool fact 2, in Lean}

\bigskip

\begin{theorem}
If $\gcd(k,n) = 1$, then
\[
  v^{2^{\,k}} = v \quad\Longleftrightarrow\quad v \in \{0,1\}.
\]
\end{theorem}

\begin{lstlisting}
lemma frob_fixed_iff_GF2 (hn : Fintype.card F = 2 ^ n)
    (hcop : Nat.Coprime k n) :
    v ^ (2 ^ k) = v ↔ (v = 0 ∨ v = 1)
\end{lstlisting}

\bigskip

\note{Coprimality $=$ minimality. The dynamical reading of a number-theory side
condition.}

\newpage


\section{The big idea}

\bigskip

The field proof of coprimality facts splits by the parity of $n$.

\bigskip

The dynamical reading does not.

\bigskip

\begin{center}
\emph{``A minimal rotation fixes only constant sequences.''}
\end{center}

\bigskip

One sentence. Every $n$. No parity.

\bigskip

That uniformity is the payoff of the encoding.

\bigskip

\note{Move the question to where it is easy. Then read the answer back.}

\newpage


\section{The functor}

\bigskip

This is not one example. It is a \emph{map of theories}.

\bigskip

\[
  \mathcal{F} : \;
  \underbrace{\text{Frobenius orbits on } \mathbb{F}_{2^{\,n}}}_{\text{algebra}}
  \;\longrightarrow\;
  \underbrace{\text{shift orbits on } \{0,1\}^{\mathbb{Z}/n\mathbb{Z}}}_{\text{dynamics}}
\]

\bigskip

At the level of orbit theories the two sides agree.

\bigskip

\begin{lstlisting}
theorem dynamics_algebra_bridge (n : ℕ) :
    MoritaEquiv (frobeniusOrbitTheory 2 n) (shiftOrbitTheory n)
\end{lstlisting}

\bigskip

So invariants cross for free:

\begin{lstlisting}
theorem bridge_transfer (n : ℕ) (I : MoritaInvariant)
    (h : I.prop (frobeniusOrbitTheory 2 n)) :
    I.prop (shiftOrbitTheory n)
\end{lstlisting}

\newpage


\section{But the functor forgets}

\bigskip

The shift sees one operation: the \textbf{multiplicative} Frobenius.

\bigskip

It never sees the field \textbf{addition}.

\bigskip

\[
  \sigma(x) = x^{2}
  \quad\text{respects}\quad \times,
  \qquad
  \text{but } +\ \text{is invisible to } \sigma .
\]

\bigskip

So the functor is \textbf{non-faithful}:

\bigskip

different additive phenomena collapse to the \emph{same} shift dynamics.

\bigskip

\note{A conjugacy of the $\times$-structure is not a conjugacy of the field.}

\newpage


\section{Why this matters}

\bigskip

APN is an \textbf{additive} property:
\[
  f(x+a) - f(x) \quad\text{is at most two-to-one.}
\]

\bigskip

The shift forgot $+$.

\bigskip

So a \emph{naive} topological conjugacy with the shift
\textbf{cannot} compute the deep APN facts.

\bigskip

The $+\,/\,\times$ clash breaks the dream of ``just work on the model.''

\bigskip

\note{The bridge is honest about its own limit.}

\newpage


\section{What survives: the spectral layer}

\bigskip

Replace the conjugacy by the \emph{transfer operator} --- the character sum.

\bigskip

\[
  W_f(a,b) \;=\; \sum_{x} (-1)^{\,\mathrm{Tr}(a x + b f(x))} .
\]

\bigskip

\begin{center}
\begin{tabular}{ll}
partition function & Walsh sum $W_f$\\[4pt]
pressure / moments & $\sum W^{2}$,\ \ $\sum W^{4}$\\[4pt]
expansion rate & $2$-adic valuation of $W$\\[4pt]
fundamental domain & Frobenius orbit (cyclotomic coset)\\
\end{tabular}
\end{center}

\bigskip

\note{Spectral data crosses the $+\,/\,\times$ clash. A bare conjugacy does not.}

\newpage


\section{APN, as a dynamical fibre condition}

\bigskip

Even without conjugacy, APN \emph{is} a statement about fibres of a dynamics.

\bigskip

\[
  \text{APN}(x \mapsto x^{d})
  \;\Longleftrightarrow\;
  \text{every fibre of } x \mapsto (x+a)^{d} + x^{d} \text{ has size} \le 2 .
\]

\begin{lstlisting}
theorem apn_small_fibers (d : ℕ) :
    IsAPN (powerFunction d : F → F) ↔
    ∀ a : F, a ≠ 0 → ∀ b : F,
      Fintype.card { x // differenceDynamics d a x = b } ≤ 2
\end{lstlisting}

\bigskip

\note{The dynamical picture frames the question --- the spectral layer answers it.}

\newpage


\section{What we claim}

\bigskip

\begin{itemize}
  \item One element $\;\mapsto\;$ one periodic sequence.\\[6pt]
  \item Frobenius $\;=\;$ shift; \ $x^{2^{k}} = \sigma^{k}$.\\[6pt]
  \item $\#\mathrm{Fix} = 2^{\gcd(k,n)}$, both sides.\\[6pt]
  \item Coprime $\;=\;$ minimal rotation $\;=\;$ fixes only constants.\\[6pt]
  \item A genuine functor (Morita bridge of orbit theories).\\[6pt]
  \item \textbf{Non-faithful}: it keeps $\times$, forgets $+$.\\[6pt]
  \item So deep APN needs the spectral layer, not a bare conjugacy.\\[6pt]
  \item All boxed statements are machine-checked.
\end{itemize}

\vfill

\noindent\textit{Read slowly. The shift was there all along.}

\end{document}
